% Full remaining Durrett Section 3.3 proof batch.
\section{Characteristic-function criteria and the moment problem}
\label{sec:annex-p15-characteristic-function-criteria-and-moments}

The source is the held author presentation
\cite[pp.~131--143]{DurrettPTEAuthor2019}.
Additional constructions below are identified where the source
states an exercise or gives no proof.

\begin{annexitem}{annex:p15.criteria-real-integration}{thm:p15-criteria-real-integration}{The full interval approximation and substitution premises}
\begin{proof}
Let \(E\) be Borel with \(\lambda(E)<\infty\).
The interval-cover construction of Lebesgue outer measure,
already used in the measure-extension theorem, gives a
countable interval cover with total length below
\(\lambda(E)+\eta\). Enlarge each interval to a bounded
open interval, choosing the total extra length below \(\eta\).
Their union \(O\) contains \(E\), and
\(\lambda(O\setminus E)<2\eta\).
Finite subunions increase to \(O\); because \(\lambda(O)<\infty\),
some finite subunion \(V\) satisfies
\(\lambda(O\setminus V)<\eta\).
Then \(\lambda(E\mathbin{\triangle}V)<3\eta\).
Choosing \(\eta\) arbitrarily small proves the first assertion.

For integrable complex \(g\), first truncate its two components
and its spatial support to make a bounded function supported
in \([-M,M]\), with arbitrarily small \(L^1\) error.
Dominated convergence justifies these truncations.
Quantizing the bounded components into a finite grid gives
a finite simple function, with error bounded by the grid
mesh times the finite interval length and a fixed component
factor. Its nonzero level sets lie in the bounded support
and have finite measure; the zero coefficient needs no set approximation.
Apply the preceding interval approximation to each set,
choosing the finitely many errors so that their sum weighted
by the absolute coefficients is arbitrarily small.
Sorting the finitely many interval endpoints expresses the
result as a step function on disjoint bounded intervals.
To approximate a step \(c\mathbf1_{(a,b)}\) continuously,
replace it by a function equal to \(c\) on
\([a+\delta,b-\delta]\), zero outside \((a,b)\), and linear
on the two end pieces, where \(0<2\delta<b-a\).
Its \(L^1\) difference from the step is \(|c|\delta\).
Choose these widths so their finite weighted sum is small.
This proves both claimed density statements, including
complex or negative coefficients.

For an interval and \(t\ne0\),
\[
 \left|\int_a^b e^{itx}dx\right|\leq2/|t|.
\]
Every finite step function therefore has Fourier integral
tending to zero as \(|t|\to\infty\).
The Fourier integrals of \(g\) and any approximating step
differ by at most their \(L^1\) distance, independently of
\(t\). First let \(|t|\to\infty\), then let this distance
tend to zero. This is the source step-function proof of
Riemann--Lebesgue, with its approximation premise supplied.

For substitution, the inverse of the stated strictly
monotone onto map is continuous: values of \(T\) at
points on either side of an inverse image enclose a
neighborhood of the image. Define on Borel subsets of \(J\)
\[
 \nu(B)=\int_I\mathbf1_B(T(x))|T'(x)|\,dx.
\]
This is a measure by monotone convergence.
For every bounded interval \((a,b]\) with endpoints in \(J\),
its inverse image is the corresponding interval between
\(T^{-1}(a)\) and \(T^{-1}(b)\), with possible endpoint
reversal. The fundamental theorem of calculus gives
\(\nu((a,b])=b-a\).
Endpoint choices do not affect the integral.
On a countable exhaustion of \(J\) by bounded intervals
these measures are finite, so interval uniqueness on each
piece gives \(\nu=\lambda|_J\); limits extend this equality
to all of \(J\).
The integral formula follows first for indicators, then
simple functions and all nonnegative functions by monotone
convergence. Applying it to \(|h|\) and the real components
proves the complex integrable version.
\end{proof}
\end{annexitem}

\begin{annexitem}{annex:p15.criteria-weak-integration}{thm:p15-criteria-weak-integration}{The source Fatou, tail-control and finite-support arguments}
\begin{proof}
Use the already proved quantile realization of weak
convergence, with \(Y_n\to Y\) almost surely and the required
marginal laws. For continuous nonnegative \(g\),
\(g(Y_n)\to g(Y)\), and Fatou gives the claimed integral
inequality, with infinite values permitted.

Under the second set of hypotheses this gives
\(\int g\,d\mu\leq C\).
For any \(\eta>0\), choose \(R\) such that
\(|h(x)|\leq\eta g(x)\) outside \([-R,R]\).
Inside that interval \(h\) is bounded.
Thus \(h\) is integrable for every law in question.
Let \(\chi_R\) be continuous, equal to one on \([-R,R]\),
zero outside \([-R-1,R+1]\), and linear on the intervening
pieces. The function \(\chi_Rh\) is bounded and continuous,
so its integrals converge. The remaining integrals have
absolute values at most \(\eta C\) for both \(\mu_n\)
and \(\mu\). Hence the upper limit of the total difference
is at most \(2\eta C\); let \(\eta\downarrow0\).

For the finite approximation, the exact quantile identity
gives
\[
 F_m(x)=\frac1m\#\{j:(j-\tfrac12)/m\leq F(x)\}.
\]
For every \(v\in[0,1]\), the fraction of these midpoint
grid levels at most \(v\) differs from \(v\) by at most
\(1/(2m)\), including equality at a grid level.
This proves the uniform bound and weak convergence.
All quantiles are finite because their levels lie in
\((0,1)\). If \(F(0)=0\), right continuity gives, at
each positive level \(u\), an \(\varepsilon>0\) with
\(F(\varepsilon)<u\); its quantile is at least
\(\varepsilon>0\). Thus the approximating atoms remain
in the required positive half-line.
\end{proof}
\end{annexitem}

\begin{annexitem}{annex:p15.criteria-complex-log}{thm:p15-criteria-complex-log}{A continuous branch and control of its integer ambiguity}
\begin{proof}
The defining series and its differentiated series converge
uniformly on each disk \(|z|\leq r<1\), by geometric
comparison. In particular
\[
 L(1+z)=\int_0^1\frac{z}{1+sz}\,ds.
\]
For fixed \(z\), differentiation along the segment \(sz\)
gives \(dL(1+sz)/ds=z/(1+sz)\).
The derivative of \(e^{L(1+sz)}/(1+sz)\) is therefore
zero, and its value at zero is one. Thus \(e^{L(1+z)}=1+z\).
The series tail gives
\[
 |L(1+z)-z|\leq\frac{|z|^2}{2(1-|z|)}\leq|z|^2
 \quad(|z|\leq1/2),
\]
and \(|L(1+z)|\leq\sum_{k\geq1}|z|^k\leq2|z|\).
The branch also equals
\(\log|1+z|+i\arctan(\operatorname{Im}(1+z)/
                         \operatorname{Re}(1+z))\):
the real part of \(1+sz\) is positive along the segment,
and direct differentiation gives the same derivative
and initial value for this expression.

For bounded \(c_n\), eventually \(|c_n/n|\leq1/2\), and
\[
 nL(1+c_n/n)=c_n+O(|c_n|^2/n)\longrightarrow c.
\]
Exponentiation proves the complex-power limit.
Finally \(e^w=1\) implies \(\operatorname{Re}w=0\)
and \(\operatorname{Im}w\in2\pi\mathbb Z\), by the
real exponential and trigonometric formulas.
A continuous function on an interval with values in
this discrete set is constant, by the intermediate
value theorem for its imaginary part.
This proves the logarithmic-lift assertion.
\end{proof}
\end{annexitem}

\begin{annexitem}{annex:p15.criteria-tight-transforms}{thm:p15-criteria-tight-transforms}{The source truncation bound and finite-grid argument}
\begin{proof}
For a required tolerance \(\varepsilon>0\), tightness gives
\(M\geq1\) with
\(\sup_i\mathbb P(|X_i|>M)<\varepsilon/4\).
The preceding characteristic-function estimate gives
\[
 |\varphi_i(t+h)-\varphi_i(t)|
 \leq\mathbb E\min\{|hX_i|,2\}
 \leq |h|M+\varepsilon/2.
\]
Take \(|h|<\varepsilon/(2M)\), proving equicontinuity.
Under weak convergence the sequence, together with its
limiting law, is tight by the already proved tightness
theorem. On a compact interval choose a finite grid
fine enough for the equicontinuity tolerance.
Pointwise convergence of transforms holds at every
grid point; eventually all finitely many errors are
small. For any other point, the difference is bounded
by the two equicontinuity errors and the nearest grid
error. This proves uniform convergence on the compact
interval. For the last assertion, the transforms are
\(e^{it/n},1\), whose difference has modulus two at
\(t=n\pi\) and never exceeds two.
\end{proof}
\end{annexitem}

\begin{annexitem}{annex:p15.criteria-uniform-sums}{thm:p15-criteria-uniform-sums}{Absolute inversion and the full polynomial construction}
\begin{proof}
The sum's transform is \((\sin t/t)^n\).
Its modulus is at most one near zero and at most
\(|t|^{-n}\) for \(|t|\geq1\), so it is integrable
for \(n\geq2\). The inversion theorem gives a bounded
continuous density. The transform is real and even;
its odd imaginary integral cancels, giving the displayed
cosine formula.

To identify the polynomial explicitly, let \(V_i=(U_i+1)/2\).
Their joint law is uniform Lebesgue measure on \((0,1)^n\).
For \(s\geq0\), the positive simplex
\(\{v_i\geq0:\sum_i v_i\leq s\}\) has volume \(s^n/n!\).
This follows by induction from Tonelli:
integrate the \((n-1)\)-dimensional volume
\((s-v_n)^{n-1}/(n-1)!\) over \(0\leq v_n\leq s\).
For \(s<0\) the volume is zero.
Inclusion--exclusion of the constraints \(v_i>1\),
translating each selected coordinate by one, therefore gives
\[
 \mathbb P\!\left(\sum_iV_i\leq s\right)
   =\frac1{n!}\sum_{j=0}^n(-1)^j\binom nj(s-j)_+^n.
\]
The translations preserve product Lebesgue measure by
the established one-coordinate affine substitutions;
the boundaries have measure zero.
Differentiating this finite expression, valid for
\(n\geq2\), gives its continuous density
\[
 \frac1{(n-1)!}\sum_{j=0}^n(-1)^j\binom nj(s-j)_+^{n-1}.
\]
It is nonnegative because the CDF is nondecreasing,
and integrates to one by its endpoint limits.
Scaling the sum by \(x=2s-n\) gives the stated
\(2^{-n}\) formula. It equals the inverse density
everywhere, since both are continuous densities of
the same law. The finite sum is polynomial away from
its integer knots \(-n+2j\). The support of the sum
is \([-n,n]\), so continuity makes the density zero
outside it, as asserted.
\end{proof}
\end{annexitem}

\begin{annexitem}{annex:p15.criteria-limit-operations}{thm:p15-criteria-limit-operations}{All three source limit operations}
\begin{proof}
In the normal case, weak convergence gives
\(e^{-v_nt^2/2}\to\varphi(t)\), where \(\varphi\) is
continuous and \(\varphi(0)=1\).
Choose \(t_0>0\) with \(\varphi(t_0)>1/2\);
this is possible since all the approximating values
are real and nonnegative. Taking real logarithms at
this single point gives
\[
 v_n\longrightarrow v=-2\log\varphi(t_0)/t_0^2
                              \in[0,\infty).
\]
At every \(t\), the limit transform is then
\(e^{-vt^2/2}\), and uniqueness identifies the
limiting normal law, including \(v=0\).

For independent rows the product rule gives
\(\varphi_{X_n+Y_n}(t)=\varphi_{X_n}(t)\varphi_{Y_n}(t)\).
The factors converge to the limiting factors, whose
product is the transform of the specified independent
limiting sum. It is continuous at zero; the continuity
theorem gives the desired weak convergence.
Only independence within each row was used.

For the almost-sure series limit, \(e^{itS_n}\to e^{itS_\infty}\)
almost surely and has modulus one. Bounded convergence
and finite independence give
\[
 \varphi_{S_\infty}(t)=\lim_n\varphi_{S_n}(t)
                      =\lim_n\prod_{j=1}^n\varphi_{X_j}(t).
\]
This proves existence of the indicated infinite product
without any nonvanishing assumption on its factors or limit.
\end{proof}
\end{annexitem}

\begin{annexitem}{annex:p15.criteria-series-examples}{thm:p15-criteria-series-examples}{Binary grids, Cantor intervals and the nonzero product}
\begin{proof}
Take versions of the digits valued in their indicated
finite sets everywhere; changing their countable union
of null exceptional sets preserves every joint law.
All the displayed series then converge absolutely at
every outcome by their deterministic geometric bounds.
For the binary series put \(B_j=(\varepsilon_j+1)/2\) and
\(V=\sum_{j\geq1}B_j2^{-j}\).
Its partial sum \(V_m\) is uniform on
\(\{k/2^m:0\leq k<2^m\}\), by independence of the \(m\)
digits, and \(0\leq V-V_m\leq2^{-m}\).
For \(0\leq x\leq1\),
\[
 \mathbb P(V_m\leq x-2^{-m})
       \leq\mathbb P(V\leq x)\leq\mathbb P(V_m\leq x).
\]
Counting the grid points makes both bounds tend to \(x\)
at \(0<x<1\); the two endpoint atom probabilities are
zero since the events of all zero or all one digits
have probabilities at most \(2^{-m}\) for every \(m\).
Thus \(V\) is uniform on \((0,1)\).
Since \(\sum_j\varepsilon_j2^{-j}=2V-1\), the product
formula from the preceding theorem and the known uniform
transform prove the asserted sinc product, including its zeros.

For the ternary construction, conditioning on the first
\(m\) digits places \(C\) in an interval
\([c,c+3^{-m}]\), where \(c\) is the corresponding partial
sum. These \(2^m\) closed intervals are pairwise disjoint:
the gap created at the first different digit exceeds the
remaining tail, leaving a positive intervening interval.
Each pattern has probability \(2^{-m}\).
Any point belongs to at most one such interval, so its
probability is at most \(2^{-m}\) for every \(m\).
The law has no atoms.
The intersection of these finite closed unions is Borel
and carries probability one, while its Lebesgue measure
is at most \((2/3)^m\) for every \(m\), hence zero.

The transform of a finite ternary sum is
\[
 \prod_{j=1}^m\frac{1+e^{2it/3^j}}2
       =\exp\!\left(it\sum_{j=1}^m3^{-j}\right)
                              \prod_{j=1}^m\cos(t/3^j).
\]
The almost-sure product theorem yields the infinite formula.
Let \(D=\prod_{j\geq1}\cos(\pi/3^j)\).
All factors are positive. For \(0\leq x\leq\pi/3\),
\(\tan x\leq4x\), since its derivative is at most four.
Thus \(0\leq-\log\cos x\leq2x^2\).
The sum of these logarithms at \(\pi/3^j\) is finite,
so \(D>0\).
At \(t=3^k\pi\), the first \(k\) cosine factors each
equal \(-1\), and the remaining product is \(D\).
The phase is \(e^{i3^k\pi/2}=i(-1)^k\).
Their product is \(iD\), independent of \(k\), proving
the source's nondecay assertion.
\end{proof}
\end{annexitem}

\begin{annexitem}{annex:p15.criteria-moment-identities}{thm:p15-criteria-moment-identities}{The source Gaussian derivative and Polya-density arguments}
\begin{proof}
Every Gaussian absolute moment is finite: on a bounded
interval its moment integrand is bounded, and outside
it the polynomial factor times \(e^{-x^2/2}\) is bounded
by a constant times \(e^{-x^2/4}\).
For example this follows because \(x^m e^{-x^2/4}\)
has a finite maximum on the positive half-line, by
differentiation.
The moment-derivative theorem therefore applies at
every order. The series
\[
 e^{-t^2/2}=\sum_{j=0}^\infty
                         \frac{(-1)^j t^{2j}}{2^j j!}
\]
and each finite-order differentiated series converge
uniformly on compact \(t\)-intervals by the ratio test.
Termwise differentiation follows by integrating the
uniformly converging derivative series, starting at zero.
The derivative of order \(2n\) at zero is
\((-1)^n(2n)!/(2^n n!)\). Comparing with
\(i^{2n}\mathbb EZ^{2n}\) gives the formula; odd
moments vanish by symmetry. The factorial cancellation
equals the product of the positive odd integers, with
the stated empty-product convention at \(n=0\).

The already normalized even Polya density gives
\[
 \int_0^\infty\frac{1-\cos u}{u^2}\,du=\pi/2.
\]
For \(x\ne0\), the affine substitution \(u=|x|t\)
therefore gives
\[
 \int_0^\infty\frac{1-\cos(tx)}{t^2}\,dt=\frac\pi2|x|;
\]
for \(x=0\) both sides are zero.
The integrands are nonnegative and jointly Borel on
\((0,\infty)\times\mathbb R\). Tonelli against Lebesgue
measure and the probability law is valid even when
the resulting integral is infinite.
Its inner expectation is \(1-\operatorname{Re}\varphi_X(t)\).
This proves the exact absolute-moment identity.
\end{proof}
\end{annexitem}

\begin{annexitem}{annex:p15.criteria-local-zero}{thm:p15-criteria-local-zero}{The local transform criterion with an explicit positive kernel}
\begin{proof}
Necessity follows from weak convergence and the
characteristic-function continuity theorem.
For sufficiency fix \(0<u\leq\delta\).
The preceding chapter's exact kernel identity and
bounded convergence give
\[
 A_n=\frac1u\int_{-u}^u(1-\operatorname{Re}\varphi_{Y_n}(t))dt
   =2\int(1-\operatorname{sinc}(ux))\,\mu_n(dx)\longrightarrow0.
\]
For every nonzero \(v\), \(\operatorname{sinc}(v)<1\):
the integral of \(\cos s\) over a nondegenerate segment
has strictly smaller absolute value than the segment's
length. The function \(2(1-\operatorname{sinc}v)\)
is continuous and tends to two as \(|v|\to\infty\).
For a fixed \(\varepsilon>0\), its infimum over
\(|v|\geq u\varepsilon\) is therefore a positive number
\(c_{u,\varepsilon}\): use a compact annulus and then
the bound at infinity.
Consequently
\[
 \mathbb P(|Y_n|\geq\varepsilon)
                 \leq A_n/c_{u,\varepsilon}\longrightarrow0.
\]
This proves convergence in probability to the constant,
and hence weak convergence, on precisely the asserted
local characteristic-function hypothesis.
\end{proof}
\end{annexitem}

\begin{annexitem}{annex:p15.criteria-independent-series}{thm:p15-criteria-independent-series}{The exact first-crossing inequality and both completeness steps}
\begin{proof}
First prove the source maximal inequality. Split the
event \(\max_{m<j\leq n}|S_{m,j}|>2a\) by its first
crossing index \(j\). Each resulting event \(A_j\)
depends only on \(X_1,\ldots,X_j\), and is independent
of the future sum \(S_{j,n}\).
On \(A_j\cap\{|S_{j,n}|\leq a\}\), the terminal sum
has \(|S_{m,n}|>a\). The \(A_j\) are disjoint, so
summing their probabilities multiplied by the minimum
future probability proves the displayed inequality.
It remains true if that minimum is zero; no division
has been used.

Now suppose \(S_n\Rightarrow\mu\), with transform
\(\varphi_n\to\varphi\) uniformly on compact intervals.
Continuity and \(\varphi(0)=1\) give a \(\delta>0\)
with \(|\varphi(t)|\geq1/2\) for \(|t|\leq\delta\).
Eventually \(\varphi_n\) is nonzero throughout this
interval. Independence gives, for \(r>s\),
\[
 \varphi_{S_r-S_s}(t)=\varphi_r(t)/\varphi_s(t).
\]
The right side tends to one uniformly there as
\(r,s\to\infty\). The positive-kernel proof of the
preceding theorem therefore gives, uniformly for
large \(r,s\),
\(\mathbb P(|S_r-S_s|>\varepsilon)\to0\).
Thus the partial sums are Cauchy in probability.

Here is the required completeness argument. For any
probability-Cauchy sequence \(Z_n\), choose increasing
indices \(n_j\) so that
\[
 \mathbb P(|Z_{n_{j+1}}-Z_{n_j}|>2^{-j})\leq2^{-j}.
\]
Borel--Cantelli makes the displayed deviation events occur
only finitely often almost surely; the summable deterministic
increment bounds then make \(Z_{n_j}\) a real Cauchy
sequence almost surely. Define its finite limit \(Z\)
on that measurable event, and zero elsewhere.
The limit is measurable by the usual limsup construction.
The probability-Cauchy condition, the triangle inequality
through \(Z_{n_j}\), and the almost-sure subsequential
limit give \(Z_n\to Z\) in probability.
Apply this to \(S_n\).

For completeness, independence upgrades that convergence
to almost sure convergence of the entire partial-sum
sequence. Put \(a_j=2^{-j}\).
Choose increasing \(N_j\) with
\[
 \sup_{r,s\geq N_j}
 \mathbb P(|S_r-S_s|>a_j)\leq2^{-j-2}.
\]
For every \(M>N_j\), the maximal inequality has its
minimum future probability at least \(1-2^{-j-2}\geq1/2\),
so
\[
 \mathbb P\!\left(\max_{N_j<k\leq M}|S_k-S_{N_j}|>2a_j\right)
                                  \leq2^{-j-1}.
\]
Let \(M\to\infty\), using continuity from below, and
then apply Borel--Cantelli in \(j\).
Almost surely, for all sufficiently large \(j\),
every \(k\geq N_j\) has \(|S_k-S_{N_j}|\leq2a_j\).
Two indices beyond \(N_j\) therefore differ by at
most \(4a_j\). The whole sequence is almost surely
Cauchy and has a finite measurable limit.
Uniqueness of limits in probability makes it the same
limit \(Z\); weak-limit uniqueness gives its law \(\mu\).
This proves every implication used in the source exercise.
\end{proof}
\end{annexitem}

\begin{annexitem}{annex:p15.criteria-pitman-law}{thm:p15-criteria-pitman-law}{Both directions with a controlled continuous logarithm}
\begin{proof}
If \(\varphi'(0)=ia\), then for each fixed real \(t\),
\(c_n=n(\varphi(t/n)-1)\to iat\).
The complex-power lemma gives
\[
 \varphi_{S_n/n}(t)=\varphi(t/n)^n=(1+c_n/n)^n\to e^{iat}.
\]
The limit is the transform of \(\delta_a\).
The continuity theorem and the constant-limit probability
bridge prove \(S_n/n\to a\) in probability.

Conversely assume that probability convergence.
Its laws converge weakly to \(\delta_a\), so their transforms
\(\Psi_n(t)=\varphi(t/n)^n\) converge uniformly to \(e^{iat}\)
on \([0,1]\), by the compact uniform theorem.
Continuity of \(\varphi\) at zero also gives
\(\sup_{0\leq t\leq1}|\varphi(t/n)-1|\to0\).
Thus for all sufficiently large \(n\) both
\(\varphi(t/n)\) and \(e^{-iat}\Psi_n(t)\) lie within
distance \(1/2\) of one throughout this interval.
The continuous functions
\[
 A_n(t)=nL(\varphi(t/n))-iat,\qquad
 D_n(t)=L(e^{-iat}\Psi_n(t))
\]
have the same exponential and are zero at zero.
The logarithmic-lift lemma gives \(A_n=D_n\), with
\(\sup_{[0,1]}|D_n|\to0\).
This equality removes, rather than assumes away, all
possible integer multiples of \(2\pi i\).

For \(h\downarrow0\), put \(n=\lfloor1/h\rfloor\)
and \(t_n=nh\). Then \(n\to\infty\), \(t_n\in(0,1]\)
and \(t_n\to1\). The preceding identity gives
\[
 \frac{L(\varphi(h))}{h}
          =ia+\frac{D_n(t_n)}{t_n}\longrightarrow ia.
\]
In particular this logarithm is \(O(h)\).
The exponential remainder
\(|e^z-1-z|\leq e^{|z|}|z|^2/2\), proved by its integral
Taylor formula, then shows
\((\varphi(h)-1)/h\to ia\).
For negative \(h\), conjugation symmetry gives the same
limit, since \(-\overline{ia}=ia\).
This proves the two-sided necessity and the source
pointwise transform assertion with its corrected sign.

For the stated example, the normalizing series is positive
and finite by comparison with \((\log2)^{-1}\sum k^{-2}\).
But \(\mathbb E|X|=c\sum_{k\geq2}(k\log k)^{-1}=\infty\),
as comparison with \(\int_2^\infty dx/(x\log x)\) shows.
For \(x\geq2\),
\[
 \mathbb P(|X|>x)
 \leq\frac{c}{\log x}\sum_{k>x}k^{-2}
 \leq\frac{2c}{x\log x}.
\]
The truncated means are
\(c\sum_{k=2}^n(-1)^k/(k\log k)\).
Their even and odd partial sums form bounded monotone
subsequences with difference tending to zero, because
the positive terms decrease to zero. Hence they converge
to the displayed finite \(a\).
The already proved full weak-tail criterion gives
\(S_n/n\to a\) in probability. The necessity just proved
then gives \(\varphi'(0)=ia\), despite the infinite
absolute first moment.
\end{proof}
\end{annexitem}

\begin{annexitem}{annex:p15.criteria-convex-derivative}{thm:p15-criteria-convex-derivative}{Secants, right continuity and the sigma-finite Stieltjes measure}
\begin{proof}
Convexity makes forward secant slopes at \(t\) nondecreasing
as their interval length increases. Their limit as that
length decreases to zero is therefore their infimum.
A backward secant to a fixed point below \(t\) bounds
this infimum below, and nonincrease of \(f\) bounds it
above by zero. Thus \(d(t)\) exists and is finite.
Convex secant comparisons give, for \(s<t\),
\[
 d(s)\leq\frac{f(t)-f(s)}{t-s}\leq d(t).
\]
In particular \(d\) is nondecreasing and bounded on
every compact subinterval of \((0,\infty)\).
For a partition \(a=x_0<\cdots<x_m=b\), multiply these
inequalities by the interval lengths and sum.
The lower and upper sums for \(d\) enclose \(f(b)-f(a)\);
their difference is at most the mesh times \(d(b)-d(a)\).
Let the mesh tend to zero to obtain the integral formula.

The right limit \(\ell=\lim_{s\downarrow t}d(s)\) is at
least \(d(t)\). The integral formula implies that every
forward secant average is at least \(\ell\).
Let its length tend to zero; then \(d(t)\geq\ell\).
This proves right continuity.
Its increasing limit at infinity is at most zero.
If that limit were negative, \(d(s)\leq-\varepsilon\)
on the entire tail for some \(\varepsilon>0\); the
integral formula would eventually force \(f\) below zero.
Thus the limit is zero.

Here is the measure construction without a finite derivative
assumption at the origin. For each integer \(j\), clip
\(d\) to its two endpoint values on
\((2^j,2^{j+1}]\).
The resulting nondecreasing right-continuous function,
minus \(d(2^j)\), is a finite Stieltjes cumulative function
of total mass \(d(2^{j+1})-d(2^j)\).
If this mass is positive, normalization gives a proper
CDF and the already proved inverse construction supplies
its measure; for mass zero use the zero measure.
It is supported on \((2^j,2^{j+1}]\).
Sum these measures over \(j\in\mathbb Z\).
The pieces are disjoint and have finite measure.
Every interval \((a,b]\) with \(0<a<b<\infty\) meets
only finitely many pieces; its increments telescope
to \(d(b)-d(a)\).
This gives the stated sigma-finite \(\eta\).
Intervals beginning at zero are defined by increasing
unions and may have infinite measure.
\end{proof}
\end{annexitem}

\begin{annexitem}{annex:p15.criteria-polya}{thm:p15-criteria-polya}{The complete source representation, approximation and direct construction}
\begin{proof}
Apply the convex-derivative lemma to \(f=\varphi|_{(0,\infty)}\).
Define \(\nu(B)=\int_B s\,\eta(ds)\).
This is a measure by monotone convergence and is sigma-finite,
because it has finite mass on each dyadic interval.
For every nonnegative measurable \(h\),
\(\int h\,d\nu=\int hs\,d\eta\): prove this for indicators
and simple functions, then pass to increasing limits.
This is the exact weighted-measure premise used in the source.
Since \(d(b)\to0\),
\[
 -d(s)=\eta((s,\infty))
                =\int_{(s,\infty)}r^{-1}\,\nu(dr).
\]
The lower endpoint is excluded, as required for the right
derivative convention.
The integral formula and the zero limit of \(\varphi\) at
infinity give, for \(t>0\),
\[
 \begin{aligned}
 \varphi(t)
 &=\int_t^\infty(-d(s))\,ds\\
 &=\int_t^\infty\int_{(s,\infty)}r^{-1}\,\nu(dr)\,ds\\
 &=\int_{(t,\infty)}(1-t/r)\,\nu(dr)
  =\int_0^\infty(1-t/r)^+\,\nu(dr).
 \end{aligned}
\]
Tonelli is justified by nonnegativity and the established
sigma-finiteness; the resulting value is finite.
Let \(t\downarrow0\). The integrands increase to one
on \((0,\infty)\), and the source continuity assumption
makes the left side tend to one. Thus \(\nu\) is a
probability. Evenness proves the formula for all real \(t\).
In the finite-piecewise-linear case, \(d\) is a step
function. Its Stieltjes measure has exactly the atoms
\(d(s_j)-d(s_j-)\) at the positive slope changes.
Weighting by \(s_j\) gives the stated finite mixture.

Retain the source approximation route.
Use the finite positive-atom approximations \(\nu_m\)
already constructed. For fixed \(t\ne0\), the function
\((1-|t|/s)^+\) on \(s>0\), extended as zero for \(s\leq0\),
is bounded and continuous on \(\mathbb R\).
For \(t=0\) use the constant extension one.
Therefore the finite mixtures
\[
 \varphi_m(t)=\int(1-|t|/s)^+\,\nu_m(ds)
\]
converge pointwise to \(\varphi(t)\).
Each is a finite mixture of characteristic functions:
if \(Y\) has the Polya density from the preceding
chapter, \(Y/s\) has transform \((1-|t|/s)^+\).
The limiting function is continuous at zero by
hypothesis. The continuity theorem supplies the law.

There is separately a direct construction.
On the product of the Polya law and \(\nu\), take
independent coordinates \(Y,S\), where \(S>0\).
The real variable \(Y/S\) is measurable.
Fubini for its bounded exponential gives exactly the
displayed mixture as its transform. This additional
construction does not replace the source approximation proof.

If the limiting value at infinity is \(c<1\), apply
the proved criterion to \((\varphi-c)/(1-c)\), and
mix its law with \(\delta_0\) using weights \(1-c,c\).
The first law has transform tending to zero, so it
has no atoms by the squared-atom identity; the mass
at zero is exactly \(c\).
For \(c=1\), monotonicity forces \(\varphi\equiv1\).
Finally \(\mathbf1_{\{0\}}\) satisfies all the other
zero-tail conditions and fails continuity at zero;
every characteristic function is continuous there,
so it supplies the stated hypothesis failure.
\end{proof}
\end{annexitem}

\begin{annexitem}{annex:p15.criteria-binomial-series}{thm:p15-criteria-binomial-series}{All real exponents and the normalized positive coefficients}
\begin{proof}
Write \(a_n=(-1)^n\binom\beta n\), with \(a_0=1\).
If the coefficients terminate, the following argument
is a finite polynomial argument. Otherwise their
successive absolute ratio tends to one.
Thus on every \(|z|\leq r<1\), the series and each
finite-order differentiated series converge uniformly
by geometric comparison.
Their coefficients satisfy
\((n+1)a_{n+1}=(n-\beta)a_n\).
The series \(H(z)=\sum a_nz^n\), for real \(|z|<1\),
therefore satisfies
\((1-z)H'(z)=-\beta H(z)\), with \(H(0)=1\).
The derivative of \((1-z)^{-\beta}H(z)\) is zero.
This proves \(H(z)=(1-z)^\beta\) on that interval,
including all real \(\beta\).

For \(0<\beta<1\), the product defining
\((-1)^{n+1}\binom\beta n\) has positive sign and
nonzero magnitude. Call it \(c_n\).
For \(0\leq z<1\),
\(\sum c_nz^n=1-(1-z)^\beta\).
Monotone convergence as \(z\uparrow1\) gives
\(\sum c_n=1\). Hence the series converges absolutely
and uniformly on \([-1,1]\).
Continuity extends the identity to both endpoints,
including the negative endpoint where positivity
of the powers themselves is not available.
\end{proof}
\end{annexitem}

\begin{annexitem}{annex:p15.criteria-stable-functions}{thm:p15-criteria-stable-functions}{The separate Polya and Spitzer source arguments}
\begin{proof}
For \(0<\alpha\leq1\), the function \(e^{-t^\alpha}\)
on \(t>0\) is decreasing, and its second derivative is
\[
 e^{-t^\alpha}
 \{\alpha^2t^{2\alpha-2}+\alpha(1-\alpha)t^{\alpha-2}\}
 \geq0.
\]
The required endpoint limits are one and zero.
Its even extension therefore satisfies Polya's criterion.
This proves the first source route.

For the separate route let \(0<\alpha<2\) and
\(\beta=\alpha/2\).
The preceding lemma gives positive weights \(c_j\)
summing to one with
\[
 \psi(t)=1-(1-\cos t)^\beta
                         =\sum_{j=1}^\infty c_j(\cos t)^j.
\]
Each \((\cos t)^j\) is the transform of a sum of \(j\)
independent fair signs, and the countable-mixture
theorem makes \(\psi\) a characteristic function.
For each fixed \(t\ne0\),
\[
 1-\cos(\sqrt2\,t\,n^{-1/\alpha})
                       \sim t^2 n^{-2/\alpha}.
\]
Consequently
\(n\{1-\psi(\sqrt2\,t\,n^{-1/\alpha})\}\to|t|^\alpha\).
The scalar product limit gives the displayed limit
of \(n\)-th powers; at \(t=0\) every value is one.
Each power is a characteristic function by scaling
and finite independent addition. Its pointwise limit
is continuous at zero, so the continuity theorem
proves existence of the limiting law.
At \(\alpha=2\), use the already computed centered
normal transform of variance two; at \(\alpha=1\)
use the Cauchy identification.
For \(\alpha>2\), the preceding chapter proved that
the zero-law consequence rules out this transform.
Finally the transform of the scaled sum is
\[
 \{e^{-|t/n^{1/\alpha}|^\alpha}\}^n=e^{-|t|^\alpha},
\]
and uniqueness proves the scaling identity.
\end{proof}
\end{annexitem}

\begin{annexitem}{annex:p15.criteria-periodic-example}{thm:p15-criteria-periodic-example}{A complete Fourier identity and normalization}
\begin{proof}
For \(0<r<1\), the logarithm series gives
\[
 F_r(x)=\sum_{k=1}^\infty\frac{r^k\sin(kx)}k
                         =-\operatorname{Im}L(1-re^{ix}).
\]
The real part of \(1-re^{ix}\) is positive, so the
branch formula in the logarithm lemma gives
\(|F_r(x)|\leq\pi/2\).
For \(0<x<2\pi\),
\[
 1-e^{ix}=2\sin(x/2)e^{i(x/2-\pi/2)},
\]
and hence \(F_r(x)\to(\pi-x)/2\) as \(r\uparrow1\).
Integrate over \([0,x]\), using uniform convergence
of the series for fixed \(r\), then bounded convergence
as \(r\uparrow1\).
The series after integration is dominated by
\(\sum 2/k^2<\infty\), so
\[
 H(x):=\sum_{k=1}^\infty\frac{1-\cos(kx)}{k^2}
                  =\frac{\pi x}{2}-\frac{x^2}{4}
                            \quad(0\leq x\leq2\pi).
\]
At \(x=\pi\) this gives
\(\sum_{k\geq1,\ k\ {\rm odd}}k^{-2}=\pi^2/8\).
For \(0\leq x\leq\pi\), subtracting the even terms yields
\[
 \sum_{k\ {\rm odd}>0}\frac{1-\cos(kx)}{k^2}
              =H(x)-\tfrac14H(2x)=\pi x/4.
\]
Thus
\[
 \frac12+\frac4{\pi^2}
       \sum_{k\ {\rm odd}>0}\frac{\cos(k\pi t)}{k^2}
                         =1-|t|\quad(|t|\leq1).
\]
The series is uniformly absolutely convergent, even
and of period two, so this equality proves the entire
periodic extension.
Pairing positive and negative odd indices gives
exactly the source's complex series.
The coefficients are nonnegative; their total mass
is \(1/2+(4/\pi^2)(\pi^2/8)=1\).
They therefore define the stated discrete probability
law and the series is its characteristic function.
Its mass \(1/2\) at zero distinguishes it from the
nonatomic Polya density law, while their transforms
agree on the whole indicated interval.
\end{proof}
\end{annexitem}

\begin{annexitem}{annex:p15.criteria-equality-sets}{thm:p15-criteria-equality-sets}{An explicit construction for every closed symmetric set}
\begin{proof}
Necessity follows from continuity of the two transforms,
their value one at zero, and their conjugation symmetry.
For the converse let \(A\) have these properties.
Define
\[
 \chi(s)=\begin{cases}s^3/(1+s^3),&s>0,\\0,&s\leq0.\end{cases}
\]
It is \(C^2\), with \(0\leq\chi\leq1\),
\(|\chi'|\leq3\), \(|\chi''|\leq24\).
For \(0\leq s\leq1\), the sharper bounds
\(\chi(s)\leq s^3\), \(|\chi'(s)|\leq3s^2\),
\(|\chi''(s)|\leq24s\) follow by differentiating
the rational expression.

Each component of this open subset of the real line is
a maximal open interval. On a component \((a,b)\) of
\((0,\infty)\setminus A\),
put \(v(t)=\chi(t-a)\chi(b-t)\).
If \(b=\infty\), omit the second factor.
Set \(v=0\) on \(A\cap[0,\infty)\).
The left endpoint is finite and belongs to \(A\),
since \(0\in A\). A finite right endpoint also belongs
to \(A\), by closedness.
Inside components the global bounds are
\[
 0\leq v\leq1,\qquad |v'|\leq6,\qquad |v''|\leq66.
\]
Near a point of \(A\), put \(d=\operatorname{dist}(t,A)\).
The factor corresponding to a nearest endpoint gives,
for \(d\leq1\),
\[
 |v(t)|\leq d^3,\quad |v'(t)|\leq6d^2,\quad
                         |v''(t)|\leq66d.
\]
These bounds show directly that the extension is
twice differentiable at every point of \(A\), with
both derivatives zero there, and that its second
derivative is continuous there.
Thus \(v\) is \(C^2\) on the positive half-line and
tends to zero at zero. It is positive exactly off \(A\).

For \(t>0\), define
\[
 \varphi_2(t)=e^{-t}(1+v(t)/160),
 \qquad \varphi_2(0)=1,
\]
and extend evenly.
Its derivative is strictly negative since
\(\varphi_2'(t)=e^{-t}[-1+(v'-v)/160]<0\).
Its second derivative is positive since
\[
 \varphi_2''(t)=e^{-t}[1+(v''-2v'+v)/160]
                     \geq e^{-t}(1-79/160)>0.
\]
It is nonnegative, tends to one at zero and to zero
at infinity. Polya's theorem therefore makes it a
characteristic function.
By construction it equals \(e^{-|t|}\) exactly on \(A\),
including the cases \(A=\mathbb R\) and \(A=\{0\}\).
This is a direct construction of the full exercise,
not just an example for one interval.
\end{proof}
\end{annexitem}

\begin{annexitem}{annex:p15.criteria-cancellation}{thm:p15-criteria-cancellation}{The countermodel and the local nonvanishing argument}
\begin{proof}
Construct \(X,Y,Z\) with the three stated laws on their
finite product probability space, making them mutually
independent. The transform of \(X\) is zero outside
\([-1,1]\), and the transforms of \(Y,Z\) agree inside
that interval. Their products with the transform of
\(X\) therefore agree everywhere.
Uniqueness gives equal laws for the two sums, whereas
\(Y,Z\) differ by the atom at zero just computed.

Conversely, equality of the laws of \(X+Y\) and \(X\)
gives \(\varphi_X(t)\varphi_Y(t)=\varphi_X(t)\).
Since \(\varphi_X(0)=1\) and it is continuous, it is
nonzero throughout some neighborhood of zero.
Division is valid there and gives \(\varphi_Y(t)=1\)
on that neighborhood. In particular it is
\(1+o(t^2)\) at zero. The previously proved zero-law
criterion forces \(Y=0\) almost surely.
\end{proof}
\end{annexitem}

\begin{annexitem}{annex:p15.criteria-complex-gamma}{thm:p15-criteria-complex-gamma}{A real-parameter proof with the branch fixed}
\begin{proof}
First put \(q=1+ib\) and
\(F(b)=\int_0^\infty t^{p-1}e^{-t}e^{-ibt}dt\).
The integral is absolutely convergent.
Its derivative integrand has modulus \(t^pe^{-t}\),
whose integral is \(\Gamma(p+1)<\infty\).
The differentiating-under-the-integral theorem gives
\[
 F'(b)=-i\int_0^\infty t^p e^{-(1+ib)t}dt.
\]
Integrate the derivative of \(t^pe^{-(1+ib)t}\)
on \([\varepsilon,R]\), then let
\(\varepsilon\downarrow0,R\to\infty\).
Both differentiated terms are absolutely integrable
and both boundary terms vanish, so
\[
 \int_0^\infty t^pe^{-(1+ib)t}dt=\frac{p}{1+ib}F(b).
\]
For \(H(b)=\tfrac12\log(1+b^2)+i\arctan b\),
one has \(H'(b)=i/(1+ib)\).
Thus \(d(e^{pH(b)}F(b))/db=0\).
Since \(F(0)=\Gamma(p)\), this proves the formula
for \(q=1+ib\) on the stated branch.
For \(q=a+ib\), \(a>0\), substitute \(u=at\).
The result is \(a^{-p}F(b/a)=\Gamma(p)q^{-p}\).
The modulus integrals are bounded by
\(\int_0^\infty t^{p-1}e^{-at}dt<\infty\).
\end{proof}
\end{annexitem}

\begin{annexitem}{annex:p15.criteria-lognormal-family}{thm:p15-criteria-lognormal-family}{All source perturbations and discrete laws}
\begin{proof}
The substitution \(x=e^s\) identifies \(f_0\) as
the pushforward density of a standard normal.
Completion of the square gives, for real \(r\),
\[
 \int_0^\infty x^rf_0(x)dx
        =\int_{\mathbb R}e^{rs}\phi(s)ds=e^{r^2/2}.
\]
For integer \(r\geq0\), the same substitution and
then translation by \(r\) yield
\[
 \int_0^\infty x^rf_0(x)\sin(2\pi\log x)dx
   =e^{r^2/2}\int_{\mathbb R}\phi(u)\sin(2\pi(u+r))du=0.
\]
The last integrand is odd.
Every perturbation factor lies in \([0,2]\).
The case \(r=0\) proves normalization; all other
integer cases prove equality of finite moments.
For distinct \(a,b\), the density difference has
a fixed nonzero sign on \(1<x<e^{1/4}\), so
their laws differ.

For the countable perturbation, the sine series
converges uniformly and absolutely with bound one.
At moment order \(r\), the sum of the absolute
integrals is at most \(e^{r^2/2}\sum|a_k|<\infty\).
It may therefore be integrated termwise.
The identity
\(\sin(k\pi(u+r))=(-1)^{kr}\sin(k\pi u)\)
makes each correction zero, proving normalization
and all the common moments.

For \(a>0\),
\[
 a^{-k}e^{-k^2/2}
       =e^{(\log a)^2/2}e^{-(k+\log a)^2/2}.
\]
For sufficiently large \(|k|\), the last factor
is at most \(e^{-k^2/8}\), a summable sequence.
Thus \(0<c_a<\infty\).
For integer \(r\geq0\), nonnegative summation and
the integer reindexing give
\[
 \begin{aligned}
 \mathbb EY_a^r
 &=c_a^{-1}\sum_{k\in\mathbb Z}a^{r-k}e^{rk-k^2/2}\\
 &=e^{r^2/2}c_a^{-1}\sum_{k\in\mathbb Z}
               a^{-(k-r)}e^{-(k-r)^2/2}=e^{r^2/2}.
 \end{aligned}
\]
If supports for \(a,b\in[1,e)\) meet,
\(\log(b/a)\) is an integer in \((-1,1)\), hence zero.
For \(b=ae^j\), reindexing gives
\(c_b=a^j e^{j^2/2}c_a\) and the same normalized
masses at the same support points. This proves
all discrete-family qualifications.
\end{proof}
\end{annexitem}

\begin{annexitem}{annex:p15.criteria-real-counterexample}{thm:p15-criteria-real-counterexample}{The whole real-line gamma cancellation}
\begin{proof}
For every nonnegative integer \(n\), substitution
\(u=x^\lambda\) gives
\[
 \int_0^\infty x^n e^{-x^\lambda}dx
                  =\lambda^{-1}\Gamma((n+1)/\lambda).
\]
The stated \(c_\lambda\) normalizes the even base
density. Every perturbed absolute moment is finite
because its factor lies between zero and two.
At even powers, including power zero, the
perturbation is odd and integrates to zero.
For odd \(n\), its full-line integral is twice
the positive-half-line sine integral.
Set \(p=(n+1)/\lambda\) and
\(\theta=\lambda\pi/2\in(0,\pi/2)\).
The complex-gamma formula gives
\[
 \int_0^\infty x^n e^{-x^\lambda}e^{-i\beta x^\lambda}dx
                   =\frac{\Gamma(p)}{\lambda(1+i\beta)^p}.
\]
On the specified branch
\(1+i\beta=(\cos\theta)^{-1}e^{i\theta}\);
the argument of its \(p\)-th power is
\((n+1)\pi/2\), an integer multiple of \(\pi\).
The displayed integral is real, so its sine
part vanishes. All moments therefore equal the
base moments, with the stated values.
Distinct parameters give different densities
on a sufficiently small positive interval
where \(0<\beta x^\lambda<\pi/2\).
\end{proof}
\end{annexitem}

\begin{annexitem}{annex:p15.criteria-compact-moments}{thm:p15-criteria-compact-moments}{The essential bound and the polynomial proof}
\begin{proof}
Put \(W=|X|\). If \(\Lambda<\infty\), then
\(\mathbb P(W<\Lambda+\varepsilon)=1\) for all
\(\varepsilon>0\); a countable decreasing
intersection gives \(W\leq\Lambda\) almost surely.
Hence \(\nu_k^{1/k}\leq\Lambda\).
For \(0<a<\Lambda\), the definition and monotonicity
of the distribution give \(p_a=\mathbb P(W\geq a)>0\).
Thus \(\nu_k^{1/k}\geq a p_a^{1/k}\to a\).
Let \(a\uparrow\Lambda\), or through all positive
finite values if \(\Lambda=\infty\).
For \(\Lambda=0\) the upper bound already proves
the result. Infinite moments cause no exception
to these nonnegative inequalities.

Suppose \(\mu\) is supported on \([-M,M]\) and
\(\eta\) has its moments.
For each \(\varepsilon>0\) and positive integer \(k\),
\[
 \eta(\{|x|>M+\varepsilon\})
 \leq m_{2k}/(M+\varepsilon)^{2k}
 \leq(M/(M+\varepsilon))^{2k}.
\]
Letting \(k\to\infty\), then using a countable
sequence \(\varepsilon\downarrow0\), forces
\(\eta\) onto the same interval.
For \(M=0\), both laws are point mass at zero.
For \(M>0\), the proved Bernstein approximation
after affine scaling to \([0,1]\) approximates
every continuous function on \([-M,M]\)
uniformly by polynomials.
Equal moments give equal polynomial integrals;
uniform approximation gives equal continuous
integrals. Apply this to restrictions of
bounded continuous real functions, and use
the continuous-ramp CDF characterization
to conclude equality of the laws.
\end{proof}
\end{annexitem}

\begin{annexitem}{annex:p15.criteria-analytic-determinacy}{thm:p15-criteria-analytic-determinacy}{The source analytic propagation with r = 0 included}
\begin{proof}
If no law has the moments, the at-most-one
assertion is immediate. If an even positive-order
moment is zero, every candidate is \(\delta_0\).
Otherwise take a candidate and put
\(\nu_j=\int|x|^j\,\mu(dx)\).
Cauchy--Schwarz gives
\(\nu_{2k+1}^2\leq m_{2k}m_{2k+2}\).
For every \(b>r\), eventually
\(\nu_j\leq(bj)^j\).
To see the odd case explicitly, first take
\(r<a<b\); the even bounds give, at odd \(j\),
\(\nu_j\leq a^j(j+1)^j\leq(bj)^j\) eventually.
The even indices give the reverse limsup inequality,
so the full absolute-moment root limsup is also \(r\).

Multiply the Taylor remainder of order \(n-1\)
by \(e^{i\theta X}\), then take expectations:
\[
 \left|\varphi(\theta+t)-
   \sum_{j=0}^{n-1}\frac{t^j}{j!}\varphi^{(j)}(\theta)\right|
                               \leq |t|^n\nu_n/n!.
\]
The exponential series gives \(n^n/n!\leq e^n\).
The remainder is therefore at most \((|t|eb)^n\)
for large \(n\), uniformly in \(\theta\).
The same bounds give absolute convergence of
the Taylor series.
For \(r>0\) and \(|t|<1/(er)\), choose
\(r<b<1/(e|t|)\); \(t=0\) is immediate.
For \(r=0\), choose such a positive \(b\)
for any fixed finite nonzero \(t\).
This proves both radius assertions.

For two candidate laws all characteristic
derivatives at zero agree.
If \(r=0\), the expansion at zero gives
equality on the entire line immediately.
If \(r>0\), put \(d=1/(3r)<1/(er)\), using \(e<3\).
The expansions at zero give equality on \([-d,d]\).
If equality holds on \([-kd,kd]\), derivatives
of every order agree on its interior and,
by continuity, at its two endpoints.
The expansions at those endpoints extend
equality to \([-(k+1)d,(k+1)d]\).
Induction covers \(\mathbb R\).
Characteristic-function uniqueness then
proves equality of the laws.
\end{proof}
\end{annexitem}

\begin{annexitem}{annex:p15.criteria-quasianalytic}{thm:p15-criteria-quasianalytic}{The full derivative-envelope argument}
\begin{proof}
Put \(a_n=M_{n-1}/M_n>0\).
The arithmetic--geometric inequality, obtained
from \(\log u\leq u-1\), applied to
\(a_1,2a_2,\ldots,na_n\) gives
\[
 (a_1\cdots a_n)^{1/n}
       \leq\frac{e}{n^2}\sum_{j=1}^n ja_j,
\]
where the factorial bound
\((n!)^{1/n}\geq n/e\) was used.
Sum from \(n=1\) to \(N\) and interchange
the finite sums. Since
\(\sum_{n=j}^\infty n^{-2}\leq j^{-2}+j^{-1}\),
the result is at most \(2e\sum_{j=1}^N a_j\).
The product on the left is \(M_0/M_n\).
Thus the assumed divergence implies
\(\sum_{n\geq1}a_n=\infty\).

Define
\[
 B(x)=\sup_{n\geq0}\frac{|f^{(n)}(x)|}{e^nM_n}.
\]
It lies in \([0,1]\) and is continuous:
the finite maxima are continuous and the
remaining terms are uniformly at most \(e^{-N}\).
Flatness gives \(B(0)=0\).
Suppose \(B(x_1)>0\) for some \(x_1>0\).
Choose \(N_0\) with \(e^{-N_0}<B(x_1)\).
For \(N\geq N_0\), let \(x_N\) be the first
point in \([0,x_1]\) where \(B=e^{-N}\).
Let \(y_N\) be the last point at or before
\(x_N\) where \(B=e^{-(N+1)}\).
Continuity gives these points and
\[
 x_{N+1}\leq y_N<x_N,\qquad
 e^{-(N+1)}<B(x)<e^{-N}\quad(y_N<x<x_N).
\]
These intervals have disjoint interiors,
with total length at most \(x_1\).

On \((y_N,x_N)\), the supremum defining \(B\)
is the maximum over \(0\leq j\leq N\),
since all higher terms are at most \(e^{-(N+1)}\).
Log-convexity makes \(M_{j+1}/M_j\)
nondecreasing. For \(y_N<p<q<x_N\), the
fundamental theorem and the modulus inequality give
\[
 \frac{|f^{(j)}(q)|}{e^jM_j}
 \leq\frac{|f^{(j)}(p)|}{e^jM_j}
       +e\frac{M_{N+1}}{M_N}\int_p^q B(s)ds
                                    \quad(j\leq N).
\]
Taking the finite maximum yields
\(B(q)\leq B(p)+L_N\int_p^q B\), where
\(L_N=eM_{N+1}/M_N\).
For \(H(q)=B(p)+L_N\int_p^q B\), one has
\(B\leq H\) and \(H'\leq L_NH\).
Differentiating \(e^{-L_N(q-p)}H(q)\) therefore gives
\[
 B(q)\leq B(p)e^{L_N(q-p)}.
\]
Let \(p\downarrow y_N,q\uparrow x_N\).
The endpoint ratio equals \(e\), so
\[
 x_N-y_N\geq L_N^{-1}
                   =M_N/(eM_{N+1})=a_{N+1}/e.
\]
Summing contradicts divergence of \(\sum a_n\)
and the finite available length.
Reflect the function to treat \(x_1<0\).
Thus \(B=0\) everywhere, and \(f\equiv0\).
The proof used only moduli and componentwise
differentiation, so it applies to complex \(f\).
\end{proof}
\end{annexitem}

\begin{annexitem}{annex:p15.criteria-carleman}{thm:p15-criteria-carleman}{A complete proof of the source's weaker condition}
\begin{proof}
If an even positive-order moment is zero,
all candidates are \(\delta_0\).
Otherwise let \(\mu,\eta\) have the moments,
and put \(f=\varphi_\mu-\varphi_\eta\).
Then \(f^{(n)}(0)=0\) and
\[
 |f^{(n)}(t)|
 \leq\int|x|^n\,\mu(dx)+\int|x|^n\,\eta(dx)
                      \leq2\sqrt{m_{2n}}=:M_n.
\]
Here \(M_0=2\).
Cauchy--Schwarz for adjacent even moments
gives \(M_n^2\leq M_{n-1}M_{n+1}\).
Also \((M_0/M_n)^{1/n}=m_{2n}^{-1/(2n)}\).
The proved quasi-analyticity lemma forces \(f=0\).
Inversion uniqueness proves \(\mu=\eta\).
This supplies the Carleman remark with its
entire additional proof, not an assumed
quasi-analyticity theorem.
Under the source growth bound, for any \(b>r\),
eventually \(m_{2k}^{1/(2k)}\leq2bk\).
Their reciprocals dominate the divergent
harmonic sequence \(1/(2bk)\).
This proves the stated implication of conditions.
\end{proof}
\end{annexitem}

\begin{annexitem}{annex:p15.criteria-moment-method}{thm:p15-criteria-moment-method}{Existence and every transferred moment}
\begin{proof}
Convergence of the second moments makes their
supremum finite, including the initial terms.
The coercive tightness theorem therefore
makes the laws tight.
Take any subsequential proper weak limit \(\eta\).
For each positive integer \(k\), apply
the proved weak-integration theorem with
\[
 h(x)=x^k,\qquad g(x)=1+|x|^{2k}.
\]
Its hypotheses hold: the ratio tends to zero
at infinity and the \(g\)-integrals are uniformly
bounded by convergence of the \(2k\)-th moments.
Hence \(\int x^k\,d\eta=m_k\), with all these
moments finite.
This proves existence of a law with the
limiting moment data. The analytic determinacy
theorem makes all such subsequential limits
that same law. The subsequence criterion
then gives weak convergence of the full sequence.
\end{proof}
\end{annexitem}

\begin{annexitem}{annex:p15.criteria-generalized-gamma}{thm:p15-criteria-generalized-gamma}{Normalization and exact root growth}
\begin{proof}
Integrability at zero requires and, with the
positive tail, is ensured by \(\alpha>-1\).
For \(\lambda>0\), substitute \(y=x^\lambda\),
so \(dx=\lambda^{-1}y^{1/\lambda-1}dy\).
The defining gamma integral gives the stated
normalizer and every moment.

For \(s\geq2\), let \(m=\lfloor s\rfloor\).
Splitting the gamma integral at one gives
\[
 \tfrac12(m-1)!\leq\Gamma(s)\leq2m!.
\]
For the lower bound use
\(\Gamma(m)-1/m\geq\tfrac12(m-1)!\);
for the upper bound use \(1+\Gamma(m+1)\leq2m!\).
These follow because on \(x\geq1\),
\(x^{m-1}\leq x^{s-1}\leq x^m\).
The recurrence \(\Gamma(m)=(m-1)!\) and
Stirling's logarithmic estimate give
\[
 \log\Gamma(s)=s\log s-s+O(\log s).
\]
Indeed \(s-m\in[0,1)\), so replacing \(m\)
by \(s\) changes \(s\log s-s\) by \(O(\log s)\).
With \(s=(n+\alpha+1)/\lambda\),
\[
 n^{-1}\log\nu_n
 =\lambda^{-1}\log n-\lambda^{-1}(\log\lambda+1)+o(1).
\]
Exponentiation proves the exact root asymptotic.
Its ratio to \(n\) tends to zero for \(\lambda>1\),
to \(1/e\) for \(\lambda=1\), and to infinity
for \(0<\lambda<1\).

A scaled absolute standard normal has this
form with \(\lambda=2,\alpha=0\);
a scaled gamma law has \(\lambda=1\),
with shape \(\alpha+1>0\).
The analytic moment condition therefore applies.
Nonzero affine changes preserve determinacy:
moments of their inverse transforms are finite
linear combinations of the given moments,
so a competing law would give a competing
law before the change. Degenerate normals
are point masses.
This proves the normal and gamma conclusions
without inferring indeterminacy merely from
failure of this sufficient bound.
\end{proof}
\end{annexitem}

\begin{annexitem}{annex:p15.criteria-stieltjes}{thm:p15-criteria-stieltjes}{The exact bijection and its criterion}
\begin{proof}
On the product with an independent fair sign,
form \(Y=\varepsilon\sqrt X\).
It is measurable, symmetric, and \(Y^2=X\).
Conversely, for a symmetric law of \(Y\),
symmetry divides the probability of \(|Y|\)
in any positive Borel set equally between its
positive and negative copies. The mass at
zero is unchanged. Thus, for \(X=Y^2\),
\(\varepsilon\sqrt X\) has the original law.
The operations are inverse on laws, including
their possible zero atoms.

When all moments of \(X\) are finite,
\(x^{j/2}\leq1+x^{\lceil j/2\rceil}\)
shows finiteness of every absolute moment
of \(Y\). Independence gives even moments
\(\nu_k\) and odd moments zero.
The displayed condition is exactly the
analytic Hamburger condition for \(Y\).
A competing nonnegative law with the same
moments gives a competing symmetric law
with all the same moments, hence the same
\(Y\) law and then the same squared law.
This proves uniqueness in the nonnegative domain.
For the generalized-gamma family, the
root asymptotic already proved gives
\[
 \frac{\nu_k^{1/(2k)}}{2k}
 \sim\tfrac12(e\lambda)^{-1/(2\lambda)}
                                  k^{1/(2\lambda)-1}.
\]
It is bounded when \(\lambda\geq1/2\),
as asserted. No change of the competing-law
domain is inferred from that calculation.
\end{proof}
\end{annexitem}

\begin{annexitem}{annex:p15.criteria-positive-counterexample}{thm:p15-criteria-positive-counterexample}{Every integer moment in the final source counterexample}
\begin{proof}
Substitution \(u=x^\lambda\) gives base integral
\(\Gamma(1/\lambda)/\lambda\).
The factor in braces lies in \([0,2]\),
so every absolute moment is finite.
For \(n\geq0\), put \(p=(n+1)/\lambda\)
and \(\theta=\lambda\pi\in(0,\pi/2)\).
The complex-gamma formula gives
\[
 \int_0^\infty x^n e^{-x^\lambda}e^{-i\beta x^\lambda}dx
                =\frac{\Gamma(p)}{\lambda(1+i\beta)^p}.
\]
Here \(1+i\beta=(\cos\theta)^{-1}e^{i\theta}\),
whose \(p\)-th power has argument
\((n+1)\pi\), an integer multiple of \(\pi\).
The integral is real, so every sine correction
vanishes, including the normalization at \(n=0\).
The base gamma moments give the stated
common moment sequence.
Different \(a\) give density differences of
fixed nonzero sign on a short positive
interval with \(0<\beta x^\lambda<\pi/2\);
hence their laws are distinct.
The specified values at zero create no
atom, since these are Lebesgue densities.
\end{proof}
\end{annexitem}
